%! Characteristics of Linear Functions — Unit 1, Lesson 1.0 — Lesson Plan
% Gradual-release plan (course-wide shape, 2026-09-01): warm-up → guided notes (with guided
% practice) → individual practice → debrief → close. There is no group activity; the "you do" is
% done alone. No tiers, no exit ticket; homework is in-packet.
\documentclass[10pt]{article}
\usepackage{algebra2-article}
\usepackage{algebra2-boxes}
\usepackage{graphicx}   % -article does NOT load graphicx; needed for thumbnails/screenshots
\graphicspath{{images/}}
\newcommand{\TallMath}[1]{$\displaystyle #1\rule[-1.4em]{0pt}{3.2em}$}
\newcommand{\CourseName}{Algebra 2}
\newcommand{\MeetingLength}{60 minutes}
\newcommand{\UnitNumberName}{Unit 1: Linear Functions \quad}
\newcommand{\LessonNumberName}{Lesson 1.0: Characteristics of Linear Functions}

\begin{document}
\begin{center}
    {\Large\bfseries \CourseName \\
    {\small\bfseries \UnitNumberName \LessonNumberName}}
\end{center}

\begin{tcolorbox}[colback=forestbg, colframe=forest, boxrule=0.9pt, arc=2mm,
    left=2mm, right=2mm, top=1.5mm, bottom=1.5mm]
    \textbf{Primary Objective:} Students can read the graph (or table, or rule) of a linear function
    and name its key characteristics --- \emph{domain} and \emph{range}, the \emph{$y$-intercept} and
    the \emph{$x$-intercept} (the \emph{zero}), the \emph{slope} as a \emph{constant rate of change},
    whether the line is \emph{increasing}, \emph{decreasing}, or \emph{constant}, and the intervals
    where the function is \emph{positive} or \emph{negative} --- justify each feature from the graph,
    and interpret it in a context. This is the ``read-a-graph'' toolkit the whole course reuses.
    \\[2pt]
    \textbf{Standards (2023 VA SOL):} \textbf{A2.F.2} --- investigate and analyze the characteristics
    of functions graphically. This Lesson~0 launches that toolkit on linear graphs:
    \textbf{A2.F.2a} (determine domain, range, zeros, and intercepts from a graph),
    \textbf{A2.F.2c} (intervals on which a graph is increasing, decreasing, or constant), and
    \textbf{A2.F.2f} (determine $f(x)$ from a graph and interpret $x$ and $f(x)$ in context). Slope
    as a constant rate of change is reactivated Algebra~1 prerequisite knowledge.
    \\[2pt]
    \textbf{Lesson model:} \emph{Gradual release} --- I do, we do, you do alone. A warm-up
    activates prior knowledge; \textbf{Guided Notes} build and name the vocabulary with the class
    on one worked context (Saturday's temperature line) and close with a \emph{Guided Practice}
    box (the ``we do''); an \textbf{Individual Practice} block on the same page is the ``you do''
    --- three problems, worked silently and alone; a \textbf{debrief} consolidates and surfaces
    the errors; \textbf{homework} extends the practice. \emph{There is no group activity in this
    lesson} --- today's independent work is done individually.
\end{tcolorbox}

\renewcommand{\arraystretch}{1.4}
\begin{skillbox}[Priority Ideas \& Skills]{goldbox}
\begin{tabularx}{\linewidth}{@{}X|X@{}}
\textbf{Priority Ideas \& Skills} & \textbf{Key Understandings (the ``why'')} \\[2pt]
\hline
\vspace{-6pt}
\begin{itemize}[leftmargin=1.1em, nosep, after=\vspace{2pt}]
  \item Find the \textbf{slope} from a graph, a table (watch the $x$-step), or any two points.
  \item Find the \textbf{$y$-intercept} by setting $x=0$ and the \textbf{zero} by solving $f(x)=0$.
  \item Decide \textbf{increasing / decreasing / constant} from the sign of the slope.
  \item Give the intervals where $f(x)>0$ and $f(x)<0$, and locate the switch at the zero.
  \item State \textbf{domain} and \textbf{range} --- for the rule, and for the story.
  \item Interpret every one of those features in a context.
\end{itemize}
&
\vspace{-6pt}
\begin{itemize}[leftmargin=1.1em, nosep, after=\vspace{2pt}]
  \item A function is \emph{linear} exactly when its rate of change is \emph{constant} --- that is
        why the graph is straight, and why \emph{any} two points give the same slope.
  \item Intercepts are where the graph meets an axis; the zero is the only place a non-horizontal
        line can change sign.
  \item \textbf{``Positive'' is not ``increasing.''} Sign describes \emph{where the graph sits};
        increasing/decreasing describes \emph{which way it moves}. A line can be positive and falling.
  \item Every non-vertical line has domain all reals; the range is all reals unless the line is
        horizontal. A \emph{story} restricts the domain to the inputs that make sense.
\end{itemize}
\end{tabularx}
\end{skillbox}

\begin{skillbox}[{Vocabulary, Concepts, \& Theorems}]{sky}
\renewcommand{\arraystretch}{1.3}
\begin{tabularx}{\linewidth}{@{}l X@{}}
\textbf{Linear function} & A function whose graph is a straight line; it can be written $f(x)=mx+b$ and has a \emph{constant} rate of change. \\
\textbf{Slope / rate of change} & $m=\dfrac{\Delta y}{\Delta x}=\dfrac{\text{rise}}{\text{run}}$: how much the output changes per unit change in the input. \\
\textbf{$y$-intercept} & The point $(0,b)$ where the graph crosses the $y$-axis --- the output when $x=0$ (the ``starting value'' in a story). \\
\textbf{$x$-intercept (zero)} & The point $(a,0)$ where the graph crosses the $x$-axis; the input $a$ with $f(a)=0$. \\
\textbf{Domain / range} & The set of allowed inputs ($x$) / the set of resulting outputs ($y$), read from the graph. \\
\textbf{Increasing / decreasing} & The graph rises ($m>0$) / falls ($m<0$) as you read it left to right. \\
\textbf{Constant function} & A horizontal line $f(x)=b$; slope $0$, neither increasing nor decreasing; range $\{b\}$. \\
\textbf{Positive / negative interval} & The $x$-values where the graph is \emph{above} the $x$-axis ($f(x)>0$) / \emph{below} it ($f(x)<0$). \\
\end{tabularx}
\end{skillbox}

\begin{fixedskillbox}[Lesson at a Glance --- \MeetingLength]{forestbg}
\renewcommand{\arraystretch}{1.25}
\begin{tabularx}{\linewidth}{@{}l c X X@{}}
\textbf{Phase} & \textbf{Min} & \textbf{Students} & \textbf{Teacher} \\
\hline
Warm-Up & 5 & Three Algebra~1 skills, alone then check with a neighbour & Circulate; debrief item~3 (the ``$+3$ each step'') aloud \\
Guided Notes & 20 & Fill the vocabulary and the four sections with the class, then \emph{Guided Practice} together & \emph{I do / we do}: build each name on the hook's temperature line; three \texttt{work} blocks at the board \\
Individual Practice & 15 & \emph{Alone and quiet}: three problems --- Sunday's falling line, the side-by-side, the car-wash card & Circulate; questions, cues, prompts --- \emph{not} answers. Note who to call on in the debrief \\
Debrief & 10 & Correct their own pages in a second colour & Work all three individual-practice problems on the board; land ``positive $\ne$ increasing'' \\
Close \& Homework & 10 & Start the homework page in class; preview 1.1 & Launch item~1 aloud, then circulate; name what changed today \\
\end{tabularx}
\end{fixedskillbox}

\begin{fixedskillbox}[Warm-Up --- Activate Prior Knowledge \& Spiral Review]{forestbg}
    The Warm-Up rehearses exactly the Algebra~1 skills the lesson leans on: \textbf{(1)} reading
    ordered pairs off a grid, including points \emph{on} an axis (the seed of ``intercept'');
    \textbf{(2)} evaluating $g(x)=3x-6$ and writing outputs as points, ending with ``for what input is
    $g(x)=0$?'' (the seed of ``zero''); and \textbf{(3)} continuing a table with a constant step and
    writing its rule $y=3x+5$ (the seed of ``slope'' and $mx+b$). Debrief item~3 aloud, briefly: a
    constant change in output for equal steps in input is what makes a rule \emph{linear} --- and
    that observation is the first line of the notes. \textbf{This page is set at 12pt} --- larger
    type than the rest of the packet, deliberately.
\end{fixedskillbox}

\begin{skillbox}[Hook]{forestbg}
    At 6:00 a.m.\ Saturday it was $-6\,^\circ$C, and the temperature rose a steady $2\,^\circ$C an
    hour. Ask: \textbf{``What was it at 7? At 8? At 10?''} ($-4$, $-2$, $2$.) \textbf{``When does it
    hit zero --- and why is that the hour that matters?''} ($h=3$, 9:00 a.m.; ice becomes water.)
    Land the idea before naming anything: two numbers tell the whole story --- \emph{where it
    starts} and \emph{how fast it changes} --- and one moment, the crossing at zero, splits the
    morning into ``icy'' and ``not.'' Today we give all of those a name and read each one three ways:
    off a graph, off a table, off a rule.
\end{skillbox}

\begin{skillbox}[Guided Notes --- the Lesson (20 min)]{forestbg}
\begin{multicols}{2}
\textbf{1. Linear means the rate of change is constant.} Fill the Saturday table and the row of
differences under it: $+2$, $+2$, $+2$. \emph{That} constant is the \textbf{slope}. Write
$m=\frac{\Delta y}{\Delta x}=\frac{\text{rise}}{\text{run}}=\frac{y_2-y_1}{x_2-x_1}$ and the form
$f(x)=mx+b$. Work the \texttt{work} block on the packet's graph of $f(x)=2x-4$ using $(1,-2)$ and
$(3,2)$: $m=2$. Ask why \emph{any} two points work, and answer it --- the rate is constant.

\textbf{2. Where the graph meets the axes.} Fill the two-row table: the \textbf{$y$-intercept}
$(0,b)$ comes from setting $x=0$ and is the \emph{starting value}; the \textbf{zero} $(a,0)$ comes
from solving $f(x)=0$ and is where the quantity runs out. Then the second \texttt{work} block:
$2x-4=0 \to x=2$. Check it against the same graph so the algebra and the picture agree.

\columnbreak
\textbf{3. Which way it moves --- and where it sits.} This is the section that matters. Fill the
two-column table ($m>0$ increasing, $m<0$ decreasing, $m=0$ constant $\mid$ above the axis positive,
below negative, sign switches only at a zero), then read them both off $f(x)=2x-4$. Close on the
gold caution box: Sunday's $g(h)=6-2h$ has $g(1)=4$ --- positive \emph{and} falling. Say it twice.

\textbf{4. Domain and range.} Non-horizontal line: all reals in, all reals out. Horizontal
$c(x)=b$: slope $0$, range $\{b\}$, and no zero unless $b=0$ --- in which case \emph{every} input is
one. In a story, the domain shrinks to the inputs that make sense.

\textbf{Guided Practice (the ``we do'').} $h(x)=-2x+6$, all features at once, including a
\texttt{work} block for the zero. Do this one \emph{with} the class --- it is the last thing they
see before working alone.
\end{multicols}
\end{skillbox}

\begin{skillbox}[Individual Practice --- \emph{On Your Own} (15 min, silent)]{redbox}
\textbf{Launch (1 min):} ``Same moves you just watched --- now nobody helps you. Three problems.
Keep the two questions apart: which way it moves, and where it sits. Stuck? Look up the page at
the notes, not at your neighbour.'' The block sits on the last page of the Guided Notes, so nothing
new is handed out.
\begin{multicols}{2}
\textbf{What students do.} \textbf{Problem~1} is Sunday evening, the falling line $T(h)=6-2h$:
rule, slope $-2$, decreasing, the intercept $(0,6)$ in words, the zero $h=3$ in a \texttt{work}
block, the sign intervals --- and then \textbf{the crux}: at $h=1$ it is $4^\circ$, \emph{positive}
and \emph{falling}, and a classmate's ``it can't be positive, it's decreasing'' to untangle in two
lines. \textbf{Problem~2} puts Saturday's and Sunday's lines side by side --- same zero, opposite
signs, slopes $2$ and $-2$ --- and asks why neither can change sign anywhere but $h=3$.
\textbf{Problem~3} is the model, the car-wash card $B(w)=40-8w$: slope and intercept in words, the
zero in a \texttt{work} block and its meaning, the inputs that make sense, and the concept check
--- a \$10 wash changes the steepness, not where the line starts.

\columnbreak
\textbf{What the teacher does.} Circulate the whole 15 minutes; keep it silent. Give a
\emph{question, cue, or prompt}, never an answer:
\begin{itemize}[leftmargin=1.1em, nosep]
  \item 1: ``Is $4$ a positive number? Is the \emph{next} hour warmer or colder? Can both be
        true?''
  \item 1: ``Where does the line touch the $h$-axis? What is $T$ there?''
  \item 2: ``Put your finger on $h=2$ for each line. Above or below? Now $h=4$.''
  \item 3: ``Where does the line start now? Same. So what \emph{did} change?''
  \item 3: ``Can you buy half a wash? So which $w$ are allowed?''
  \item Finished early: ``Write Sunday's rule starting from the 7~p.m.\ reading instead of the
        4~p.m.\ one. Does it agree?''
\end{itemize}
Note names as you go: one student for each of the three problems, to put work on the board in the
debrief.
\end{multicols}
\end{skillbox}

\boxguard
\begin{skillbox}[Debrief (10 min --- what goes on the board, in a second color)]{forestbg}
Students correct their own Individual Practice in a second colour as you go. Work all three, but
protect the time for items~2 and~3.
\begin{enumerate}[leftmargin=1.3em, nosep, itemsep=2pt]
  \item \textbf{Problem~1 --- the whole toolkit on a falling line.} Label slope $-2$, intercept
        $(0,6)$, zero $(3,0)$, positive before $h=3$ and negative after, directly on the graph.
  \item \textbf{The must-land moment: ``positive'' is not ``increasing.''} $4^\circ$ and
        falling. \emph{Sits} vs.\ \emph{moves}. Point back at the gold caution box in notes
        section~3. This is the misconception the unit test will probe; say it twice.
  \item \textbf{Problem~2 --- the same zero, opposite signs.} Both lines cross at $h=3$; before
        and after swap. Draw the conclusion out loud: a line changes sign only at its zero, so it
        changes sign at most once.
  \item \textbf{Problem~3(d) --- rate vs.\ starting value.} A \$10 wash does not move the
        $y$-intercept; it steepens the line, so the card empties \emph{sooner} ($w=4$). A student
        who says ``the line starts lower'' is reading a rate as a starting value --- send them to
        the graph, not to the algebra.
\end{enumerate}
If time is short, cut item~1 (they can self-check from the key) and protect items~2 and~3.
\end{skillbox}

\boxguard
\begin{skillbox}[Active Monitoring --- Watch For]{redbox}
\begin{itemize}[leftmargin=1.2em, nosep]
  \item swapping the intercepts --- reading the $y$-intercept off the $h$-axis or vice versa (Individual Practice~1);
  \item calling Sunday's rate ``$2$'' instead of ``$-2$'' --- the sign \emph{is} the direction (Individual Practice~1);
  \item fusing ``positive'' with ``increasing'' (Individual Practice~1, homework~2 and~6) --- the target misconception;
  \item writing the flat line's range as ``all reals'' or giving it a zero (notes~4, homework~4);
  \item reading homework~3's slope as $3$ --- the $x$-step is $2$, so $m=\tfrac{\Delta y}{\Delta x}=\tfrac32$;
  \item giving the story's domain when the \emph{rule}'s is asked for (notes~4, homework~5c) --- push ``forget the story; what does the rule allow?''
\end{itemize}
Cold-call prompts: ``Where is the zero, and what does $T$ equal there?'' ``Rising or falling --- and
what is the sign of the rate?'' ``Name an hour where $T<0$.''
\end{skillbox}

\boxguard
\begin{skillbox}[Reinforcement \& Extension]{goldbox}
\textbf{Homework} (in the packet, collected next class): read every feature off a \emph{rule}
($f(x)=2x+2$); a contrast pair of graphs with the same zero at $x=1$, closing with ``why can the
sign change only at the zero?''; a \emph{table} whose $x$-step is $2$ (so $m=\tfrac32$, not $3$);
the constant function $c(x)=-2$ and its boundary case $b=0$; a model, $W(t)=600+150t$, whose zero
$t=-4$ is real on the line and meaningless in the story; and an SOL-style multiple-choice item on
$y=-3x+6$. \textbf{Extension:} build the rule from a $y$-intercept and a zero, and explain why no
line can be positive only on a bounded interval. \textbf{DeltaMath override:} this content is well
covered there --- swap in a DeltaMath set on slope, intercepts, and interpreting linear graphs if you
prefer auto-graded practice. \textbf{Preview:} Lesson~1.1, \emph{Linear Functions: slope, rate of
change, and forms of a line} --- today they \emph{read} those features off a line; next they run it
backwards and \emph{write} the equation.
\end{skillbox}

% ── Teacher notes — one per component, in packet order ──
\begin{teachernote}[Warm-Up]
Items 1 and 2 plant ``a point on an axis'' and ``the input that makes the output $0$'' without
naming them; item~3 plants ``constant step $\Rightarrow$ rule $mx+b$.'' Debrief item~3 only, and
only the step --- then carry that exact observation straight into notes section~1, where it becomes
the definition. Five minutes; if item~1 runs long, skip the ``on the axis'' follow-up.
\end{teachernote}

\begin{teachernote}[Guided Notes]
Twenty minutes, paced by section: 3 for the vocabulary and hook, 4 for section~1 (including the
slope \texttt{work} block), 3 for section~2, 5 for section~3, 2 for section~4, and 3 for the
Guided Practice box worked \emph{with} the class. The
\emph{whole lesson turns on section~3}: build the two-column table so the class sees that
``moves'' and ``sits'' are different questions before the gold caution box gives them a case where
the answers disagree. Two \texttt{work} blocks are printed in the notes (slope from two points, and
$2x-4=0$) and a third in the practice box ($-2x+6=0$) --- students write in the reserved space and
the key prints the same steps there, so the two files paginate alike (two more sit in Individual
Practice). If you fall behind, compress section~4 to its first bullet --- but do \emph{not} cut
section~3's caution box; Individual Practice~1 and homework~2 both target it.
\end{teachernote}

\begin{teachernote}[Individual Practice]
Fifteen minutes, and \textbf{keep it silent} --- this is the only block all period where a student
finds out what they can do without help. Budget roughly 6, 3, and 6 minutes. Problem~1 is
deliberately the same toolkit as the notes, on a \emph{falling} line, so the sign of the slope has
to be read rather than assumed; its crux --- $4^\circ$ and falling --- is the lesson's target
misconception, and a student who writes ``negative'' there has fused \emph{sits} with \emph{moves}.
Send them to the gold caution box, not to a rule. Common slips: reading Sunday's slope as $+2$;
writing the negative interval as ``$h<3$'' with the sides swapped; in 3(c), giving $0\le w\le5$
rather than whole washes. Fast finishers rewrite Sunday's rule from the 7~p.m.\ reading and check
it agrees; that is the work you want on the board in the debrief.
\end{teachernote}

\begin{teachernote}[Homework]
Item~3 is the one that separates: the $x$-step is $2$, so students who divide the $y$-differences by
$1$ get $3$ instead of $\tfrac32$. Item~4(c) is the boundary case --- $c(x)=b$ has a zero only when
$b=0$, and then \emph{every} input is one; expect ``one zero'' as the common wrong answer. Item~5(b)
is the story-vs-math point again, and item~6 is the formative check: sort into ``chose C (got it) /
chose B (fused \emph{positive} with \emph{right of the zero}) / chose A (fused \emph{positive} with
\emph{increasing}) / chose D (read the slope as the intercept)'' to decide whether Lesson~1.1 opens
with a one-minute re-look at the Sunday line. Two \texttt{work} blocks (item~5b and extension~a).
With the group activity cut, the last 10 minutes are for \emph{starting} this page in class ---
launch item~1 aloud, then circulate; the rest is due next class.
\end{teachernote}

\end{document}
